\documentclass{beamer}

\usepackage{amsmath,amssymb,amsfonts}
\usepackage{physics}
\usepackage{bm}
\usepackage{graphicx}
\usepackage{tikz}
\usetikzlibrary{arrows.meta,positioning}

\title{Reinforcement Learning for Quantum Control}
\subtitle{A Physics-Oriented Introduction with Applications to Quantum Sensors}
\author{Morten Hjorth-Jensen}
\date{March 17, 2026}

\begin{document}

%================================================
\frame{\titlepage}

%================================================
\begin{frame}{Outline}
\begin{enumerate}
\item Why reinforcement learning?
\item Basic ingredients: states, actions, rewards, policies
\item Value functions and Bellman equations
\item Q-learning, policy gradients, actor--critic
\item Worked RL example with equations and pseudocode
\item Model-aware reinforcement learning
\item Reinforcement learning for quantum control
\item Applications to quantum sensing
\end{enumerate}
\end{frame}

%================================================
\begin{frame}{Motivation}
Many problems in physics can be phrased as optimization problems:
\begin{itemize}
\item variational calculations,
\item control of dynamical systems,
\item adaptive measurements,
\item maximizing sensitivity in quantum sensors.
\end{itemize}

However, in many realistic settings:
\begin{itemize}
\item the objective is not known analytically,
\item gradients are unavailable or noisy,
\item decisions must be taken sequentially,
\item present choices affect future opportunities.
\end{itemize}

Reinforcement learning (RL) addresses precisely this type of problem.
\end{frame}


%================================================
\begin{frame}{What is Reinforcement Learning?}
  An agent interacts with an environment over discrete time steps:
  \[
  s_0,a_0,r_1,s_1,a_1,r_2,\dots
  \]

  Core ingredients:
  \begin{itemize}
  \item \textbf{State} $s_t$: information available at time $t$,
  \item \textbf{Action} $a_t$: decision taken by the agent,
  \item \textbf{Reward} $r_{t+1}$: feedback from the environment,
  \item \textbf{Policy} $\pi(a|s)$: rule for choosing actions.
  \end{itemize}

  Goal:
  \[
  \max_\pi \; \mathbb{E}_\pi\!\left[\sum_{t=0}^{\infty}\gamma^t r_{t+1}\right]
  \]
  where $0\le \gamma <1$ is a discount factor.
\end{frame}

%================================================
\begin{frame}{Sequential Decision Making}
  In ordinary optimization we choose parameters once.

  In reinforcement learning we choose a \emph{sequence} of actions:
  \[
  a_0,a_1,a_2,\dots
  \]

  This matters because:
  \begin{itemize}
  \item actions change the future state,
  \item the future state changes which actions are useful later,
  \item rewards may be delayed.
  \end{itemize}

  Thus RL is best viewed as \textbf{optimization over trajectories or strategies}.
\end{frame}

%================================================
\begin{frame}{Markov Decision Process}
  The standard framework is a Markov Decision Process (MDP).

  Elements:
  \begin{itemize}
  \item state space $\mathcal S$,
  \item action space $\mathcal A$,
  \item transition probability $P(s'|s,a)$,
  \item reward function $R(s,a)$ or $R(s,a,s')$.
  \end{itemize}

  Markov property:
  \[
  P(s_{t+1}|s_t,a_t,s_{t-1},a_{t-1},\dots)=P(s_{t+1}|s_t,a_t)
  \]

  This is analogous to stochastic evolution equations in statistical physics.
\end{frame}

%================================================
\begin{frame}{Return and Value Functions}
  Define the return from time $t$:
  \[
  G_t = \sum_{k=0}^\infty \gamma^k r_{t+k+1}
  \]

  The \textbf{state-value function} under policy $\pi$ is
  \[
  V^\pi(s)=\mathbb{E}_\pi[G_t\mid s_t=s]
  \]

  The \textbf{action-value function} is
  \[
  Q^\pi(s,a)=\mathbb{E}_\pi[G_t\mid s_t=s,\; a_t=a]
  \]

  The \textbf{advantage} is
  \[
  A^\pi(s,a)=Q^\pi(s,a)-V^\pi(s)
  \]

  Interpretation: how much better is action $a$ than the typical action at state $s$?
\end{frame}

%================================================
\begin{frame}{Bellman Equation}
  The value function satisfies a recursive relation:
  \[
  V^\pi(s)=\sum_a \pi(a|s)\sum_{s'} P(s'|s,a)
  \left[r(s,a,s')+\gamma V^\pi(s')\right]
  \]

  This is the Bellman equation.

  Interpretation:
  \[
  \text{value} = \text{immediate reward} + \text{discounted future value}
  \]

  This is a self-consistency relation, reminiscent of mean-field equations and iterative schemes in physics.
\end{frame}

%================================================
\begin{frame}{Optimal Bellman Equation}
  Define the optimal value
  \[
  V^*(s)=\max_\pi V^\pi(s)
  \]
  and optimal action-value function
  \[
  Q^*(s,a)=\max_\pi Q^\pi(s,a).
  \]

  Then
  \[
  V^*(s)=\max_a \sum_{s'} P(s'|s,a)
  \left[r(s,a,s')+\gamma V^*(s')\right]
  \]

  and
  \[
  Q^*(s,a)=\sum_{s'} P(s'|s,a)
  \left[r(s,a,s')+\gamma \max_{a'}Q^*(s',a')\right].
  \]

  This is the basis for dynamic programming and Q-learning.
\end{frame}

%================================================
\begin{frame}{Value Iteration and Policy Iteration}
  \textbf{Value iteration:}
  \[
  V_{k+1}(s)=\max_a \sum_{s'} P(s'|s,a)\left[r+\gamma V_k(s')\right]
  \]

  \textbf{Policy iteration:}
  \begin{enumerate}
  \item Evaluate $V^\pi$ for current policy,
  \item Improve the policy using
    \[
    \pi_{\text{new}}(s)=\arg\max_a Q^\pi(s,a)
    \]
  \end{enumerate}

  These are exact if the model is known and the state-action spaces are manageable.
\end{frame}

%================================================
\begin{frame}{Q-Learning}
  When the model is unknown, we can learn directly from sampled transitions.

  Q-learning update:
  \[
  Q(s_t,a_t)\leftarrow Q(s_t,a_t) +
  \alpha\Big[r_{t+1}+\gamma \max_{a'}Q(s_{t+1},a')-Q(s_t,a_t)\Big]
  \]

  Here:
  \begin{itemize}
  \item $\alpha$ is the learning rate,
  \item the term in brackets is the temporal-difference error.
  \end{itemize}

  This is a stochastic approximation to the Bellman optimality equation.
\end{frame}

%================================================
\begin{frame}{Exploration vs Exploitation}
  A central issue in RL is the exploration--exploitation tradeoff.

  \begin{itemize}
  \item \textbf{Exploitation}: choose the best-known action,
  \item \textbf{Exploration}: try other actions to discover better strategies.
  \end{itemize}

  Common strategies:
  \begin{itemize}
  \item $\varepsilon$-greedy,
  \item softmax/Boltzmann exploration,
  \item entropy regularization.
  \end{itemize}

  This is analogous to balancing local minimization and global search in complex optimization landscapes.
\end{frame}

%================================================
\begin{frame}{Why Policy Gradient Methods?}
  Q-learning is natural for discrete action spaces.

  But in physics and control problems we often have:
  \begin{itemize}
  \item continuous control amplitudes,
  \item stochastic policies,
  \item large state spaces,
  \item differentiable parameterizations.
  \end{itemize}

  Then it is natural to parameterize the policy directly:
  \[
  \pi_\theta(a|s)
  \]
  and optimize the expected return
  \[
  J(\theta).
  \]
\end{frame}

%================================================
\begin{frame}{Policy Gradient Setup}
  Define
  \[
  J(\theta)=\mathbb{E}_{\tau\sim \pi_\theta}[R(\tau)]
  \]
  where $\tau=(s_0,a_0,s_1,a_1,\dots)$ is a trajectory and
  \[
  R(\tau)=\sum_{t=0}^\infty \gamma^t r_{t+1}.
  \]

  Trajectory probability:
  \[
  P_\theta(\tau)=\rho(s_0)\prod_{t=0}^\infty \pi_\theta(a_t|s_t)P(s_{t+1}|s_t,a_t)
  \]

  We want to compute
  \[
  \nabla_\theta J(\theta).
  \]
\end{frame}

%================================================
\begin{frame}{Policy Gradient Theorem: Step 1}
  Start with
  \[
  J(\theta)=\sum_\tau P_\theta(\tau)R(\tau).
  \]

  Differentiate:
  \[
  \nabla_\theta J(\theta)=\sum_\tau \nabla_\theta P_\theta(\tau)\,R(\tau).
  \]

  Use the log-derivative identity:
  \[
  \nabla_\theta P_\theta(\tau)=P_\theta(\tau)\nabla_\theta \log P_\theta(\tau).
  \]

  Hence
  \[
  \nabla_\theta J(\theta)=\sum_\tau P_\theta(\tau)\nabla_\theta \log P_\theta(\tau)\,R(\tau).
  \]

  Therefore
  \[
  \nabla_\theta J(\theta)=\mathbb{E}_{\tau\sim \pi_\theta}\left[\nabla_\theta \log P_\theta(\tau)\,R(\tau)\right].
  \]
\end{frame}

%================================================
\begin{frame}{Policy Gradient Theorem: Step 2}
  Now expand the log trajectory probability:
  \[
  \log P_\theta(\tau)=\log \rho(s_0)+
  \sum_{t=0}^\infty \log \pi_\theta(a_t|s_t)+
  \sum_{t=0}^\infty \log P(s_{t+1}|s_t,a_t).
  \]

  Only the policy depends on $\theta$, so
  \[
  \nabla_\theta \log P_\theta(\tau)=
  \sum_{t=0}^\infty \nabla_\theta \log \pi_\theta(a_t|s_t).
  \]

  Substituting,
  \[
  \nabla_\theta J(\theta)=
  \mathbb{E}_{\tau\sim \pi_\theta}
  \left[
    \left(
    \sum_{t=0}^\infty \nabla_\theta \log \pi_\theta(a_t|s_t)
    \right)
    R(\tau)
    \right].
  \]
\end{frame}

%================================================
\begin{frame}{Policy Gradient Theorem: Step 3}
  We can refine this by assigning to each time step only the future return from that point onward:
  \[
  G_t=\sum_{k=t}^\infty \gamma^{k-t}r_{k+1}.
  \]

  Then
  \[
  \nabla_\theta J(\theta)=
  \mathbb{E}_{\pi_\theta}
  \left[
    \sum_{t=0}^\infty \nabla_\theta \log \pi_\theta(a_t|s_t)\,G_t
    \right].
  \]

  This is the basis of the REINFORCE algorithm.

  A more compact form is
  \[
  \nabla_\theta J(\theta)=
  \mathbb{E}_{s\sim d^\pi,\,a\sim \pi_\theta}
  \left[
    \nabla_\theta \log \pi_\theta(a|s)\,Q^{\pi_\theta}(s,a)
    \right]
  \]
  where $d^\pi(s)$ is the discounted state visitation distribution.
\end{frame}

%================================================
\begin{frame}{Variance Reduction and Baselines}
  The estimator can have high variance.

  We may subtract a baseline $b(s)$ without changing the expectation:
  \[
  \nabla_\theta J(\theta)=
  \mathbb{E}\left[
    \nabla_\theta \log \pi_\theta(a|s)\left(Q^\pi(s,a)-b(s)\right)
    \right].
  \]

  A natural choice is
  \[
  b(s)=V^\pi(s),
  \]
  which gives
  \[
  \nabla_\theta J(\theta)=
  \mathbb{E}\left[
    \nabla_\theta \log \pi_\theta(a|s)\,A^\pi(s,a)
    \right].
  \]

  This motivates actor--critic methods.
\end{frame}

%================================================
\begin{frame}{REINFORCE Algorithm}
  \textbf{Basic Monte Carlo policy gradient algorithm}

  For each episode:
  \begin{enumerate}
  \item Generate a trajectory using $\pi_\theta$,
  \item Compute returns $G_t$,
  \item Update parameters:
    \[
    \theta \leftarrow \theta + \alpha
    \sum_t \nabla_\theta \log \pi_\theta(a_t|s_t)\,G_t
    \]
  \end{enumerate}

  Advantages:
  \begin{itemize}
  \item simple,
  \item direct,
  \item works with stochastic policies.
  \end{itemize}

  Disadvantage:
  \begin{itemize}
  \item large variance.
  \end{itemize}
\end{frame}

%================================================
\begin{frame}{Actor--Critic Methods}
  Actor--critic combines:
  \begin{itemize}
  \item \textbf{Actor}: policy $\pi_\theta(a|s)$,
  \item \textbf{Critic}: estimate of value function $V_\phi(s)$.
  \end{itemize}

  Temporal-difference error:
  \[
  \delta_t=r_{t+1}+\gamma V_\phi(s_{t+1})-V_\phi(s_t)
  \]

  Policy update:
  \[
  \theta \leftarrow \theta + \alpha_\theta
  \nabla_\theta \log \pi_\theta(a_t|s_t)\,\delta_t
  \]

  Critic update:
  \[
  \phi \leftarrow \phi - \alpha_\phi \nabla_\phi \delta_t^2
  \]

  This is often much more efficient than pure Monte Carlo policy gradients.
\end{frame}

%================================================
\begin{frame}{Worked RL Example: Grid World}
  Consider a particle moving on a 1D lattice:
  \[
  x \in \{0,1,2,3,4\}.
  \]

  Actions:
  \[
  a \in \{\text{left},\text{right}\}.
  \]

  Goal:
  \begin{itemize}
  \item reach $x=4$,
  \item avoid spending too many steps,
  \item possibly avoid a penalty region.
  \end{itemize}

  Reward example:
  \[
  r=
  \begin{cases}
    +10, & x=4,\\
    -1, & \text{each intermediate step}.
  \end{cases}
  \]

  This is simple enough to analyze explicitly, but already contains the core RL ideas.
\end{frame}

%================================================
\begin{frame}{Worked Example: Bellman Equations}
  If the policy is deterministic and always moves right, then
  \[
  V(4)=0
  \]
  if we stop at the goal, and for $x<4$,
  \[
  V(x)=-1+\gamma V(x+1).
  \]

  Thus
  \begin{align*}
    V(3) &= -1,\\
    V(2) &= -1+\gamma(-1),\\
    V(1) &= -1+\gamma(-1-\gamma),\\
    V(0) &= -1+\gamma(-1-\gamma-\gamma^2).
  \end{align*}

  So even this toy model already illustrates how future consequences are encoded recursively.
\end{frame}

%================================================
\begin{frame}{Worked Example: Q-Learning Equations}
  Suppose the agent experiences
  \[
  s_t=1,\qquad a_t=\text{right},\qquad r_{t+1}=-1,\qquad s_{t+1}=2.
  \]

  Then the update is
  \[
  Q(1,\text{right}) \leftarrow Q(1,\text{right}) +
  \alpha\left[-1 + \gamma \max_a Q(2,a)-Q(1,\text{right})\right].
  \]

  If initially all $Q=0$, then
  \[
  Q(1,\text{right}) \leftarrow \alpha(-1).
  \]

  Repeated experience gradually propagates information backward from the goal.
\end{frame}

%================================================
\begin{frame}{Worked Example: Pseudocode}
  \small
\begin{verbatim}
  Initialize Q(s,a) arbitrarily
For each episode:
    Initialize state s
    While s is not terminal:
        Choose action a from Q using epsilon-greedy
        Take action a, observe reward r and next state s'
        Q(s,a) <- Q(s,a) + alpha *
                  [r + gamma * max_a' Q(s',a') - Q(s,a)]
        s <- s'
\end{verbatim}

\begin{itemize}
\item This is model-free.
\item No knowledge of transition probabilities is required.
\item The environment itself generates the data.
\end{itemize}
\end{frame}

%================================================
\begin{frame}{Worked Example: Policy Gradient Pseudocode}
  \small
\begin{verbatim}
  Initialize policy parameters theta
  For each episode:
  Generate trajectory using pi_theta
  For each time step t:
  Compute return G_t
  Update:
          theta <- theta + alpha *
                 sum_t [grad_theta log pi_theta(a_t|s_t) * G_t]
\end{verbatim}

With a baseline:
\[
\theta \leftarrow \theta + \alpha
\sum_t \nabla_\theta \log \pi_\theta(a_t|s_t)\,
\big(G_t - b(s_t)\big)
\]

This becomes especially natural when actions are continuous.
\end{frame}

%================================================
\begin{frame}{Model-Free vs Model-Aware Reinforcement Learning}
  Two broad philosophies:

  \textbf{Model-free RL}
  \begin{itemize}
  \item learn directly from transitions,
  \item no explicit model of $P(s'|s,a)$ is used,
  \item examples: Q-learning, REINFORCE.
  \end{itemize}

  \textbf{Model-aware (or model-based) RL}
  \begin{itemize}
  \item use or learn a model of the dynamics,
  \item plan using that model,
  \item combine prediction and control.
  \end{itemize}

  In physics, model-aware RL is often especially attractive because we frequently know at least part of the dynamics.
\end{frame}

%================================================
\begin{frame}{What is Model-Aware Reinforcement Learning?}
  In model-aware RL we assume access to:
  \begin{itemize}
  \item known transition probabilities $P(s'|s,a)$, or
  \item an approximate learned model $\hat P_\phi(s'|s,a)$.
  \end{itemize}

  Then we can:
  \begin{itemize}
  \item simulate hypothetical futures,
  \item estimate returns before acting,
  \item plan over multiple time steps,
  \item use differentiable dynamics when available.
  \end{itemize}

  This is close in spirit to standard computational physics:
  \[
  \text{equations of motion} \;\Rightarrow\; \text{simulation} \;\Rightarrow\; \text{optimization}.
  \]
\end{frame}

%================================================
\begin{frame}{Benefits of Model-Aware RL}
  Advantages:
  \begin{itemize}
  \item better sample efficiency,
  \item ability to plan ahead,
  \item easier incorporation of prior physical knowledge,
  \item natural interface with control theory.
  \end{itemize}

  Disadvantages:
  \begin{itemize}
  \item model bias if the model is wrong,
  \item errors may accumulate during planning,
  \item learning an accurate model may be difficult.
  \end{itemize}

  Thus model-aware RL is especially useful when the dynamics are fairly well understood, but the control task remains hard.
\end{frame}

%================================================
\begin{frame}{Model-Aware RL and Physics}
  In many physics applications we know the dynamics:
  \[
  \dot x = f(x,u)
  \]
  or in quantum mechanics
  \[
  i\hbar \frac{d}{dt}\ket{\psi(t)}=H(u(t))\ket{\psi(t)}.
  \]

  Then the environment is not a black box:
  \begin{itemize}
  \item we can simulate trajectories,
  \item estimate gradients,
  \item evaluate candidate policies offline.
  \end{itemize}

  Model-aware RL thus sits naturally between:
  \begin{itemize}
  \item optimal control,
  \item stochastic optimization,
  \item machine learning.
  \end{itemize}
\end{frame}

%================================================
\begin{frame}{Transition to Quantum Systems}
  Quantum control problems are naturally sequential:
  \begin{itemize}
  \item apply pulse,
  \item observe state evolution,
  \item choose next pulse,
  \item optimize final performance.
  \end{itemize}

  This makes RL conceptually attractive.

  At the same time, quantum dynamics are often known from the Schr\"odinger equation, which makes model-aware RL particularly relevant.
\end{frame}

%================================================
\begin{frame}{Quantum Control}
  Quantum dynamics:
  \[
  i\hbar \frac{d}{dt}\ket{\psi(t)} = H(t)\ket{\psi(t)}
  \]

  Controlled Hamiltonian:
  \[
  H(t)=H_0+\sum_k u_k(t)H_k
  \]

  Typical goals:
  \begin{itemize}
  \item state preparation,
  \item gate synthesis,
  \item suppression of decoherence,
  \item enhancement of sensor performance.
  \end{itemize}
\end{frame}

%================================================
\begin{frame}{RL Formulation for Quantum Control}
  Mapping:
  \begin{itemize}
  \item \textbf{State}: wavefunction, density matrix, or measurement record,
  \item \textbf{Action}: control amplitude or pulse choice,
  \item \textbf{Reward}: fidelity, energy, Fisher information, robustness.
  \end{itemize}

  Examples of terminal rewards:
  \[
  R = |\langle \psi_{\mathrm{target}}|\psi(T)\rangle|^2
  \]
  or
  \[
  R = F_Q(T)
  \]
  for a quantum sensing task.
\end{frame}

%================================================
\begin{frame}{Quantum Sensors and Fisher Information}
  A quantum sensor estimates a parameter $\theta$.

  Precision is bounded by the quantum Cram\'er--Rao bound:
  \[
  \Delta \theta \ge \frac{1}{\sqrt{\nu F_Q}}
  \]
  where:
  \begin{itemize}
  \item $F_Q$ is the quantum Fisher information,
  \item $\nu$ is the number of repetitions.
  \end{itemize}

  Thus maximizing $F_Q$ is a natural control objective for reinforcement learning.
\end{frame}

%================================================
\begin{frame}{Worked Quantum Example: Two-Level Sensor}
  Consider a qubit sensing a magnetic field $B$:
  \[
  H(t)=\frac{\omega_0}{2}\sigma_z + \frac{B}{2}\sigma_z + u_x(t)\sigma_x.
  \]

  Interpretation:
  \begin{itemize}
  \item $\omega_0$: bare splitting,
  \item $B$: unknown signal,
  \item $u_x(t)$: control field chosen by the agent.
  \end{itemize}

  Goal:
  \begin{itemize}
  \item accumulate phase information about $B$,
  \item keep coherence as long as possible,
  \item maximize final sensitivity.
  \end{itemize}
\end{frame}

%================================================
\begin{frame}{Quantum Sensor as an RL Problem}
  Discretize time into steps $t_1,\dots,t_N$.

  Possible actions:
  \[
  a_t \in \{-u,0,+u\}.
  \]

  Possible state representation:
  \begin{itemize}
  \item Bloch vector $(\langle \sigma_x\rangle,\langle \sigma_y\rangle,\langle \sigma_z\rangle)$,
  \item time step,
  \item previous measurement outcomes.
  \end{itemize}

  Possible reward:
  \[
  R = F_Q(T)
  \]
  or
  \[
  R=\left|\frac{\partial}{\partial B}\langle \sigma_y(T)\rangle\right|.
  \]
\end{frame}

%================================================
\begin{frame}{Why RL for Quantum Sensors?}
  RL can learn how to:
  \begin{itemize}
  \item alternate between free precession and control,
  \item suppress unwanted evolution,
  \item balance phase accumulation against decoherence,
  \item adapt measurements based on partial information.
  \end{itemize}

  The learned strategy may resemble:
  \begin{itemize}
  \item Ramsey interferometry,
  \item spin echo,
  \item dynamical decoupling,
  \item adaptive phase estimation.
  \end{itemize}
\end{frame}

%================================================
\begin{frame}{Model-Aware RL for Quantum Control}
  Quantum systems are a prime example of model-aware RL.

  If the Hamiltonian is known, we can:
  \begin{itemize}
  \item simulate trajectories offline,
  \item evaluate hypothetical pulse sequences,
  \item use differentiable propagators,
  \item combine RL with standard optimal-control ideas.
  \end{itemize}

  This is especially attractive for:
  \begin{itemize}
  \item quantum control,
  \item calibration,
  \item sensor design,
  \item adaptive metrology.
  \end{itemize}
\end{frame}

%================================================
\begin{frame}{Summary}
  Reinforcement learning provides:
  \begin{itemize}
  \item a framework for sequential decision making,
  \item value-based and policy-based algorithms,
  \item a natural extension of stochastic optimization,
  \item direct relevance for quantum control and quantum sensing.
  \end{itemize}

  For physicists, the key conceptual bridge is:
  \[
  \text{variational optimization} + \text{stochastic dynamics} + \text{control}
  \;\Rightarrow\; \text{reinforcement learning}.
  \]
\end{frame}

%================================================
\begin{frame}{Suggested Reading}
  \begin{itemize}
  \item R. Sutton and A. Barto, \emph{Reinforcement Learning: An Introduction}
  \item M. Bukov et al., \emph{Reinforcement Learning in Different Phases of Quantum Control}, PRX 8, 031086 (2018)
  \item M. Cerezo et al., \emph{Variational Quantum Algorithms}, Nat. Rev. Phys. 3, 625 (2021)
  \item J. Biamonte et al., \emph{Quantum Machine Learning}, Nature 549, 195 (2017)
  \end{itemize}
\end{frame}

\end{document}
